\section{Measurability, finite covers, and a scalar benchmark}
\label{app:measurable}
\subsection{Finite-label implementation}
A finite codebook $C=\{c_1,\ldots,c_M\}$ in a metric space has a Borel
nearest-centre selector whenever each distance $s\mapsto d(s,c_j)$ is
Borel. Define its label by the least index attaining the minimum. The set
on which label $j$ is selected is the intersection of the finitely many
sets $\{d(s,c_j)<d(s,c_i)\}$ for $i<j$ and
$\{d(s,c_j)\le d(s,c_i)\}$ for $i>j$. These sets are Borel. Thus the
finite-label recursions used in this paper do not require an unproved
measurable choice of a nearest point in an infinite set.

Codebook construction and codebook execution are different questions.
Existence of finitely many representatives in nonempty cells is an ordinary
finite choice. When an algorithmic construction is claimed, one must also
supply a representation and a feasibility procedure for those cells. In
Section~\ref{sec:refreshing}, rational calibration gives a rational
polytope; finite rational linear feasibility supplies this additional
construction. This does not charge an unbounded list of history-dependent
centres to a read-only program.

\subsection{Conditional distributions and common versions}
The existence of regular conditional distributions on standard Borel
spaces is a classical input \cite{Kallenberg2021}. Its relevant conclusion
is existence of a kernel under a fixed joint law. It does not provide one
canonical version on every null history, nor does it alone give a version
continuous in a parameter or uniformly valid on an uncountable family of
null sets. The statements in the main text accordingly take one of three
forms: an explicit positive finite-report formula; a compact continuous
realization as in Theorem~\ref{thm:compact-quotient}; or an already supplied
jointly measurable predictive realization with its applicable quantifiers.

There is no hidden claim that all quotient images of a Borel space are
standard Borel. The evaluation image into a countable product can be given
its induced measurable structure, but to use ordinary kernel calculus on
it one must verify that the image and descended maps have the required
measurability. The compact theorem verifies precisely that fact in one
useful class. Keeping the full posterior remains a valid alternative when
compression of the state is not yet justified.

\subsection{Quantization after a readout change}
\begin{proposition}[Bi-Lipschitz readout comparison]\label{prop:readout-comparison}
Let $S$ be a metric space with a probability $\nu$, and let $p:S\to\R^q$
be an injective readout satisfying
\[
 a\,d(s,t)\le\norm{p(s)-p(t)}_2\le b\,d(s,t),\qquad a,b>0.
\]
If $S$ has an $M$-point cover of radius $r$ with centres in $S$, the
checkpoint prediction risk is at most $b^2r^2$. If some measurable map
$T:\R^q\to\R^k$ has Lipschitz constant $L>0$, and the law of $Tp(s)$
has density at most $A$, that risk is at least
\[
 \frac{1}{2L^2}(2Av_kM)^{-2/k}.
\]
The second assertion does not require decoder centres to belong to $p(S)$.
\end{proposition}
\begin{proof}
For the upper bound map the cover centres through $p$ and use the upper
Lipschitz estimate. For any decoder centres $z_j\in\R^q$, the inequality
$\norm{p(s)-z_j}_2\ge L^{-1}\norm{Tp(s)-Tz_j}_2$ reduces the lower
bound to Lemma~\ref{lem:small-ball}. Take the infimum over the original
centres. The constant $a$ is useful for comparing intrinsic covers in both
directions, but it is not needed in this particular density-based lower
bound. In particular, the argument does not require inversion of a
near-singular statistical covariance.
\end{proof}

A linear projection or bounded coordinate map can be used for $T$.
A coordinate change whose Lipschitz constant tends to infinity changes the
bound; it cannot be called a harmless normalization at a degeneracy.

\begin{proposition}[Exact scalar uniform quantization]\label{prop:uniform-quantization}
For $U$ uniform on an interval of length $a>0$, the optimal squared
$M$-centre quantization error is $a^2/(12M^2)$, for every integer $M\ge1$.
\end{proposition}
\begin{proof}
Order the centres. Their nearest-centre cells intersect the interval in
intervals, allowing empty cells. If a cell has length $l_j$, its contribution
is minimized by its midpoint and equals $l_j^3/(12a)$. Since
$\sum_j l_j=a$, convexity gives
$\sum_j l_j^3\ge a^3/M^2$. Hence the total error is at least
$a^2/(12M^2)$. Equal-length cells with their midpoints attain that value.
Arbitrary randomized encoders cannot improve it by
Theorem~\ref{thm:quantization}.
\end{proof}

This exact formula is a benchmark for the meaning of label cardinality,
not a claim that the dynamic posterior density in
Section~\ref{sec:refreshing} is uniform. The latter is bounded by its
actual Jacobian and can depend on the preceding history. Its two-sided
order, rather than an exact optimizer constant, is what we prove.

\subsection{Predictable error budgets and random horizons}
The stopped version of Lemma~\ref{lem:telescoping} is often more useful
than the deterministic sum. In a comparison with common visible stopping
rule $\sigma$, append an absorbing terminal symbol. The two conditional
kernels coincide once the symbol has appeared. Their discrepancy at time
$t$ is therefore bounded by $\1_{\{t<\sigma\}}e_t(H_t)$, a measurable
function of the common prefix. The telescoping proof integrates it against
the actual prefix law. This gives the expectation of the active-step sum,
without replacing a random horizon by an uncharged successful count.

If the active-step sum is bounded by a deterministic budget $B$ on every
path, its expectation is at most $B$. If only an expected budget is known,
that is the exact kind of budget furnished by the theorem. A tail-budget
statement requires an additional probability estimate. These are distinct
resource certificates even though the same positive summands occur in
all three.
